\clearpage
\oldpage[I]{214}
\phantomsection
\addcontentsline{toc}{section}{Bibliography}
\section*{Bibliography}

\begin{thebibliography}{22}

\bibitem[1]{I-1}
H.~\textsc{Cartan} and C.~\textsc{Chevalley},
\emph{Séminaire de l'École Normale Supérieure},
8th year (1955--56), \emph{Géométrie algébrique}.

\bibitem[2]{I-2}
H.~\textsc{Cartan} and S.~\textsc{Eilenberg},
\emph{Homological Algebra}, Princeton Math. Series (Princeton University
Press), 1956.

\bibitem[3]{I-3}
W.~L.~\textsc{Chow} and J.~\textsc{Igusa},
\emph{Cohomology theory of varieties over rings},
\emph{Proc. Nat. Acad. Sci. U.S.A.}, vol.~XLIV (1958), pp.~1244--1248.

\bibitem[4]{I-4}
R.~\textsc{Godement}, \emph{Théorie des faisceaux}, Actual. Scient. et Ind.,
no.~1252, Paris (Hermann), 1958.

\bibitem[5]{I-5}
H.~\textsc{Grauert},
\emph{Ein Theorem der analytischen Garbentheorie und die Modulräume komplexer
Strukturen}, \emph{Publ. Math. Inst. Hautes Études Scient.},
no.~5, 1960.

\bibitem[6]{I-6}
A.~\textsc{Grothendieck}, \emph{Sur quelques points d'algèbre homologique},
\emph{Tôhoku Math. Journ.}, vol.~IX (1957), pp.~119--221.

\bibitem[7]{I-7}
A.~\textsc{Grothendieck},
\emph{Cohomology theory of abstract algebraic varieties},
\emph{Proc. Intern. Congress of Math.}, pp.~103--118, Edinburgh (1958).

\bibitem[8]{I-8}
A.~\textsc{Grothendieck},
\emph{Géométrie formelle et géométrie algébrique},
\emph{Séminaire Bourbaki}, 11th year (1958--59), exposé 182.

\bibitem[9]{I-9}
M.~\textsc{Nagata},
\emph{A general theory of algebraic geometry over Dedekind domains},
\emph{Amer. Math. Journ.}: I, vol.~LXXVIII, pp.~78--116 (1956); II, vol.~LXXX,
pp.~382--420 (1958).

\bibitem[10]{I-10}
D.~G.~\textsc{Northcott}, \emph{Ideal theory}, Cambridge Univ. Press, 1953.

\bibitem[11]{I-11}
P.~\textsc{Samuel}, \emph{Commutative algebra} (Notes by D.~Herzig), Cornell
Univ., 1953.

\bibitem[12]{I-12}
P.~\textsc{Samuel}, \emph{Algèbre locale}, Mém. Sci. Math.,
no.~123, Paris, 1953.

\bibitem[13]{I-13}
P.~\textsc{Samuel} and O.~\textsc{Zariski},
\emph{Commutative algebra}, 2 vols., New York (Van Nostrand), 1958--60.

\bibitem[14]{I-14}
J.-P.~\textsc{Serre}, \emph{Faisceaux algébriques cohérents},
\emph{Ann. of Math.}, vol.~LXI (1955), pp.~197--278.

\bibitem[15]{I-15}
J.-P.~\textsc{Serre},
\emph{Sur la cohomologie des variétés algébriques},
\emph{Journ. de Math.} (9), vol.~XXXVI (1957), pp.~1--16.

\bibitem[16]{I-16}
J.-P.~\textsc{Serre},
\emph{Géométrie algébrique et géométrie analytique},
\emph{Ann. Inst. Fourier}, vol.~VI (1955--56), pp.~1--42.

\bibitem[17]{I-17}
J.-P.~\textsc{Serre},
\emph{Sur la dimension homologique des anneaux et des modules noethériens},
\emph{Proc. Intern. Symp. on Alg. Number theory}, pp.~176--189, Tokyo-Nikko,
1955.

\bibitem[18]{I-18}
A.~\textsc{Weil}, \emph{Foundations of algebraic geometry}, Amer. Math. Soc.
Coll. Publ., no.~29, 1946.

\bibitem[19]{I-19}
A.~\textsc{Weil}, \emph{Numbers of solutions of equations in finite fields},
\emph{Bull. Amer. Math. Soc.}, vol.~LV (1949), pp.~497--508.

\bibitem[20]{I-20}
O.~\textsc{Zariski},
\emph{Theory and applications of holomorphic functions on algebraic varieties
over arbitrary ground fields}, Mem. Amer. Math. Soc., no.~5
(1951).

\bibitem[21]{I-21}
O.~\textsc{Zariski}, \emph{A new proof of Hilbert's Nullstellensatz},
\emph{Bull. Amer. Math. Soc.}, vol.~LIII (1947), pp.~362--368.

\bibitem[22]{I-22}
E.~\textsc{Kähler}, \emph{Geometria Arithmetica},
\emph{Ann. di Mat.} (4), vol.~XLV (1958), pp.~1--368.

\end{thebibliography}
